% Generated draft: requires mathematical and source review.
\documentclass[11pt]{article}
\usepackage[margin=25mm]{geometry}
\usepackage{fontspec}
\setmainfont{Latin Modern Roman}
\usepackage{amsmath,amssymb,mathtools,unicode-math}
\setmathfont{Latin Modern Math}
\usepackage{enumitem,needspace,xurl,hyperref}
\setlength{\parindent}{0pt}
\setlength{\parskip}{5pt}
\setlength{\emergencystretch}{3em}
\newcommand{\N}{\mathbb{N}}
\newcommand{\Z}{\mathbb{Z}}
\newcommand{\Q}{\mathbb{Q}}
\newcommand{\R}{\mathbb{R}}
\newcommand{\C}{\mathbb{C}}
\newcommand{\F}{\mathbb{F}}
\newcommand{\A}{\mathbb{A}}
\newcommand{\PP}{\mathbb{P}}
\newcommand{\E}{\mathbb{E}}
\newcommand{\eps}{\varepsilon}
\newcommand{\dd}{\,\mathrm{d}}
\begin{document}
\section*{UnsolvedMath ID: 800002}
\textbf{Source problem code:} catalogue entry
\subsection{Definitions and mathematical statement}
A semigroup is a nonempty set with an associative binary operation. A zero is an element \(0\) satisfying \(0s=s0=0\) for every element \(s\). For this draft, a semigroup \(S\) with zero is called \(0\)-simple if \(S\ne\{0\}\), its only nonempty two-sided ideals are \(\{0\}\) and \(S\), and \(S^2\ne\{0\}\). The source does not explicitly specify this convention.

Under the usual semigroup interpretation, a semigroup \(T\) divides a semigroup \(S\), written \(T\preccurlyeq S\), if there exist a nonempty subsemigroup \(U\subseteq S\) and a surjective semigroup homomorphism \(\varphi:U\to T\). For a class \(\mathcal C\), define
\[
\operatorname{Div}(\mathcal C)=\{T:\text{there exists }S\in\mathcal C\text{ with }T\preccurlyeq S\}.
\]
Here homomorphisms are ordinary semigroup homomorphisms; preservation of a distinguished zero is not imposed separately. The source does not define ``divisor,'' so this interpretation requires confirmation.

Let \(G\) be a finite group with identity \(1_G\), and let \(I\) and \(\Lambda\) be finite nonempty sets. Adjoin an absorbing zero to \(G\). Given a matrix \(P=(p_{\lambda i})_{\lambda\in\Lambda,\,i\in I}\) over \(G\cup\{0\}\) with no row or column consisting entirely of zeros, the Rees matrix semigroup
\[
M^0(G;I,\Lambda;P)=(I\times G\times\Lambda)\cup\{0\}
\]
has absorbing zero and multiplication
\[
(i,g,\lambda)(j,h,\mu)=
\begin{cases}
(i,gp_{\lambda j}h,\mu),&p_{\lambda j}\ne0,\\
0,&p_{\lambda j}=0.
\end{cases}
\]
A finite Brandt semigroup is a semigroup isomorphic to
\[
B(G,I)=M^0(G;I,I;P),\qquad
p_{\lambda i}=\begin{cases}1_G,&\lambda=i,\\0,&\lambda\ne i.\end{cases}
\]
Let \(\mathcal B\) be the class of these finite Brandt semigroups. Let \(\mathcal R_{01}\) be the class of finite semigroups isomorphic to some \(M^0(G;I,\Lambda;P)\) as above for which
\[
p_{\lambda i}\in\{0,1_G\}\qquad(\lambda\in\Lambda,\ i\in I).
\]
Thus the word ``ones'' is interpreted as group identity entries, and membership requires the existence of such a Rees representation. The source does not restrict the finite group to be trivial or specify a preferred representation.

\textbf{Problem (S. Kublanovsky).} Describe, up to isomorphism, each of the two classes
\[
\operatorname{Div}(\mathcal B)
\qquad\text{and}\qquad
\operatorname{Div}(\mathcal R_{01}).
\]
Equivalently, for each class give necessary and sufficient conditions on a finite semigroup \(T\) for \(T\) to be a homomorphic image of a subsemigroup of some member of that class. The divisor \(T\) is not required to be \(0\)-simple or to have a zero.

The source adds that divisors of the class of all finite \(0\)-simple semigroups are described by positive universal formulas, referring to \([\mathrm{Kubl}]\). This is retained as supplied background; neither those formulas nor a complete reference for \([\mathrm{Kubl}]\) are supplied. The request does not explicitly require the two desired descriptions to use positive universal formulas. The record labels the problem open; this is a supplied status, not an independently verified assessment.
\subsection{Short English statement}
Describe exactly which finite semigroups arise as homomorphic images of subsemigroups of finite Brandt semigroups, and which arise in the same way from finite \(0\)-simple semigroups admitting a Rees matrix with entries only zero and the group identity. The problem is attributed to S. Kublanovsky. The divisor convention, the meaning of the matrix entries, and the intended form of the descriptions require source confirmation.
\subsection{Sources}
\begin{itemize}
\item Local UnsolvedMath registry, permanent ID \(800002\), exact source code \texttt{catalogue entry}, title \emph{The Lyapin's Notebook --- Problem 2}; supplied upstream revision \texttt{d371ee917ac9647fbb06cd6fc529deec6f3dda82}. The supplied record labels the problem \texttt{open}.
\item Bershadsky, Kublanovsky, Mashevitzky, \emph{The Lyapin's Notebook: Unsolved Problems in Semigroup Theory} (2026), Problem 2, attributed within the statement to S. Kublanovsky. Supplied source URL: \url{https://arxiv.org/abs/2604.04763}. The background records access on 2026-10-16 and extraction by \texttt{arxiv\_latex}; these are supplied provenance notes, not a claim of independent source access or verification. The registry lists Bershadsky, Kublanovsky, Mashevitzky in its proposer field and 2026 in its proposed-year field.
\item Reference label \([\mathrm{Kubl}]\), supplied for the assertion about positive universal formulas describing divisors of all finite \(0\)-simple semigroups. No bibliographic details or separate URL for this reference are present in the supplied record.
\end{itemize}
\end{document}
